\documentclass[a4paper,12pt]{article}

\usepackage[T2A]{fontenc}
\usepackage[utf8]{inputenc}
\usepackage[russian]{babel}
\usepackage{amsmath,amssymb,amsthm,mathtools}

\usepackage{helvet}
\renewcommand{\familydefault}{\sfdefault}

\usepackage{icomma}
\usepackage[a4paper,margin=2.2cm,headheight=25pt]{geometry}
\usepackage{microtype}

\usepackage{tikz}
\usepackage{tkz-euclide}
\usepackage{pgfplots}
\pgfplotsset{compat=1.18}

\usepackage{tabularx,booktabs,longtable}
\usepackage{enumitem}
\usepackage{hyperref}
\usepackage{parskip}
\usepackage{fancyhdr}
\usepackage{setspace}
\usepackage{xcolor}

% -------------------------------------------------
% Палитра
% -------------------------------------------------
\definecolor{accent}{HTML}{008080}
\definecolor{darkaccent}{HTML}{004D4D}
\definecolor{textgray}{HTML}{333333}
\definecolor{softgray}{HTML}{EEF3F3}
\definecolor{muted}{HTML}{6D7A7A}

\hypersetup{hidelinks}

\onehalfspacing
\setlength{\parindent}{0pt}

\setlist[itemize]{
  leftmargin=1.6em,
  itemsep=0.25em,
  topsep=0.3em
}

\setlist[enumerate]{
  leftmargin=1.8em,
  itemsep=0.35em,
  topsep=0.4em
}

\renewcommand{\arraystretch}{1.25}

% -------------------------------------------------
% Колонтитулы
% -------------------------------------------------
\pagestyle{fancy}
\fancyhf{}
\lhead{\small\color{muted}ОГЭ по математике}
\rhead{\small\color{muted}София Кальней}
\cfoot{\small\thepage}

\renewcommand{\headrulewidth}{0.3pt}
\renewcommand{\headrule}{%
  \hbox to\headwidth{%
    \color{softgray}%
    \leaders\hrule height \headrulewidth\hfill
  }%
}

% -------------------------------------------------
% Служебные команды
% -------------------------------------------------
\newcommand{\topic}[1]{%
  \vspace{0.8em}%
  {\Large\bfseries\color{darkaccent}#1}\par
  \vspace{0.25em}%
  {\color{accent}\rule{\linewidth}{0.6pt}}\par
  \vspace{0.45em}%
}

\newcommand{\key}[1]{%
  \textbf{\color{darkaccent}#1}%
}

\begin{document}

% =================================================
% ТИТУЛЬНЫЙ ЛИСТ
% =================================================
\begin{titlepage}
\thispagestyle{empty}
\centering

\vspace*{2.6cm}

{\fontsize{18}{22}\selectfont
\color{accent}\bfseries
ЧЕК-ЛИСТ\par}

\vspace{0.8cm}

{\fontsize{27}{32}\selectfont
\bfseries\color{textgray}
Комплексное повторение алгебры\par}

\vspace{0.35cm}

{\Large\color{darkaccent}
уравнения, неравенства, системы,\\
текстовые задачи и графики\par}

\vspace{1.5cm}

{\large ОГЭ по математике\par}

\vspace{0.4cm}

{\large\bfseries София Кальней\par}

\vfill

{\large Лёвин Артём Александрович\par}

{\normalsize
эксклюзивно для Софии Кальней\par}

\vspace{0.35cm}

{\normalsize \today\par}

\vspace{1.1cm}

\begin{tikzpicture}[x=1cm,y=1cm]
  \draw[
    accent,
    line width=1.1pt
  ]
    (-6.3,0)
    .. controls (-4.4,0.65) and (-2.4,-0.45) ..
    (-0.3,0.08)
    .. controls (1.7,0.60) and (3.5,-0.38) ..
    (6.3,0.12);

  \draw[
    softgray,
    line width=0.7pt
  ]
    (-6.3,-0.18)
    .. controls (-4.0,0.20) and (-1.8,-0.62) ..
    (0.2,-0.18)
    .. controls (2.4,0.28) and (4.6,-0.35) ..
    (6.3,-0.12);
\end{tikzpicture}

\end{titlepage}

% =================================================
% ОСНОВНОЙ ЧЕК-ЛИСТ
% =================================================

\topic{1. Быстрый выбор метода}

Перед преобразованиями определите тип задачи и цель первого шага.

\begin{table}[ht]
\centering
\caption{Карта решений}
\begin{tabularx}{\linewidth}{@{}>{\bfseries}p{0.27\linewidth}X@{}}
\toprule
Ситуация & Первый рациональный шаг \\
\midrule

Уравнение с числовыми знаменателями
&
Умножить обе части на общий числовой знаменатель и убрать дроби.
\\

Полиномиальное неравенство
&
Перенести всё в одну сторону, получить ноль справа, разложить выражение на множители.
\\

Рациональное неравенство
&
Перенести всё в одну сторону, привести левую часть к одной дроби, указать ОДЗ, применить метод интервалов.
\\

Система уравнений
&
Выбрать сложение или вычитание так, чтобы одна переменная исчезла максимально быстро.
\\

Система неравенств
&
Решить каждое неравенство отдельно и взять пересечение множеств решений.
\\

Совокупность
&
Решить каждую часть и взять объединение множеств решений.
\\

Текстовая задача
&
Ввести искомую величину и заполнить таблицу по основной формуле темы.
\\

Кусочная функция
&
Разобрать каждую ветвь в своей области и отдельно проверить граничные точки.
\\

\bottomrule
\end{tabularx}
\end{table}

\key{Контрольная привычка.}
После каждого преобразования проверяйте ОДЗ, знаки и условия включения
граничных точек.

% -------------------------------------------------

\topic{2. Уравнения: дроби, свёртка и разложение на множители}

\key{Уравнение с числовыми дробями.}
Если знаменатели содержат только числа, удобно сразу умножить
уравнение на их НОК.

\textbf{Пример.}
\[
\frac{x^2}{4}+3x+8=0.
\]

Умножаем обе части на \(4\):
\[
x^2+12x+32=0.
\]

Разложим левую часть:
\[
x^2+12x+32=(x+4)(x+8).
\]

Следовательно,
\[
\boxed{x=-4,\;-8}.
\]

\key{Квадратный трёхчлен внутри выражения.}
Полезно сразу проверять возможность применения формул сокращённого
умножения:
\[
x^2+4x+4=(x+2)^2.
\]

Например,
\[
(x-1)(x+2)^2-4(x+2)=0.
\]

Выносим общий множитель:
\[
(x+2)\bigl((x-1)(x+2)-4\bigr)=0.
\]

Получаем две ветви:
\[
x+2=0
\]
или
\[
(x-1)(x+2)-4=0.
\]

Вторая ветвь приводит к квадратному уравнению:
\[
x^2+x-6=0.
\]

Таким образом,
\[
\boxed{x=-3,\;-2,\;2}.
\]

\key{Сокращение алгебраических дробей.}
Ограничения исходного выражения сохраняются после сокращения.

Например,
\[
\frac{(x-3)(x+1)}{x-3}>0,
\qquad x\ne3.
\]

После сокращения рассматриваем
\[
x+1>0
\]
с дополнительным условием
\[
x\ne3.
\]

Поэтому
\[
x>-1,
\qquad x\ne3.
\]

% -------------------------------------------------

\topic{3. Неравенства и метод интервалов}

\key{Алгоритм для произведения или дроби:}

\begin{enumerate}
  \item Перенесите все выражения в одну сторону, чтобы справа стоял \(0\).

  \item Разложите числитель и знаменатель на множители, если это
  упрощает анализ.

  \item Найдите нули числителя.

  \item Найдите нули знаменателя и внесите их в ОДЗ.

  \item Отметьте все критические точки на числовой прямой.

  \item Определите знак на одном интервале и восстановите остальные
  знаки с учётом кратностей корней.

  \item Выберите интервалы требуемого знака и правильно оформите
  границы.
\end{enumerate}

\textbf{Пример 1.}
\[
(x-3)(x-3-\sqrt5)<0.
\]

Критические точки:
\[
x=3,
\qquad
x=3+\sqrt5.
\]

Обе точки открытые, поскольку неравенство строгое.

Для \(x=0\):
\[
(0-3)(0-3-\sqrt5)>0.
\]

На левом интервале стоит знак \(+\).
Далее знаки чередуются:

\begin{center}
\begin{tikzpicture}[x=2cm,y=1cm]

\draw[->,line width=0.8pt]
  (-0.6,0)--(3.5,0)
  node[right] {$x$};

\draw[fill=white,line width=0.8pt]
  (0.8,0) circle (2.2pt);

\draw[fill=white,line width=0.8pt]
  (2.6,0) circle (2.2pt);

\node[below=4pt] at (0.8,0) {$3$};
\node[below=4pt] at (2.6,0) {$3+\sqrt5$};

\node[above=5pt] at (0.05,0) {$+$};
\node[above=5pt] at (1.7,0) {$-$};
\node[above=5pt] at (3.1,0) {$+$};

\draw[accent,line width=2pt]
  (0.84,-0.23)--(2.56,-0.23);

\end{tikzpicture}
\end{center}

Следовательно,
\[
\boxed{x\in(3;\,3+\sqrt5)}.
\]

\key{Чётная кратность корня.}
Если множитель встречается в чётной степени, знак при переходе через
соответствующую точку сохраняется.

Например,
\[
(x-2)^2
\]
имеет чётную кратность в точке \(x=2\).

\subsection*{Рациональное неравенство}

Рассмотрим
\[
\frac1x-\frac1{x-3}>0.
\]

Сначала выписываем ОДЗ:
\[
x\ne0,
\qquad
x\ne3.
\]

Приводим выражение к одной дроби:
\[
\frac{x-3-x}{x(x-3)}>0.
\]

Получаем
\[
\frac{-3}{x(x-3)}>0.
\]

Критические точки:
\[
0,\qquad3.
\]

Обе точки всегда открытые, поскольку обращают знаменатель в ноль.

На интервале \((0;3)\), например при \(x=1\),
\[
\frac{-3}{1(1-3)}>0.
\]

Следовательно,
\[
\boxed{x\in(0;3)}.
\]

\key{О домножении неравенств.}
Домножение на выражение с неизвестным знаком без анализа его знака
может изменить направление неравенства. Для рациональных
неравенств надёжный маршрут состоит в приведении к одной дроби и
применении метода интервалов.

% -------------------------------------------------

\topic{4. Системы и совокупности}

\subsection*{Система уравнений}

Ищите действие, которое быстро устраняет одну переменную.

\textbf{Пример.}
\[
\begin{cases}
3x^2+y=7,\\
4x^2-y=0.
\end{cases}
\]

Складываем уравнения:
\[
7x^2=7.
\]

Отсюда
\[
x^2=1,
\qquad
x=\pm1.
\]

Подставляем \(x^2=1\) во второе уравнение:
\[
4-y=0,
\]
поэтому
\[
y=4.
\]

Ответ:
\[
\boxed{(1;4),\;(-1;4)}.
\]

\subsection*{Система неравенств}

Система требует пересечения множеств решений:
\[
A\cap B.
\]

\subsection*{Совокупность}

Совокупность требует объединения:
\[
A\cup B.
\]

Для системы удобно искать участки, где штриховки совпали.
Для совокупности собираются все участки, принадлежащие хотя бы одному
множеству.

\textbf{Пример.}
\[
\begin{cases}
\displaystyle
\frac{(x+2)(x-3)}{x+5}\ge0,
\\[0.6em]
(x-2)^2(x+7)\le0.
\end{cases}
\]

Первое неравенство имеет решение
\[
(-5;-2]\cup[3;+\infty).
\]

Для второго неравенства учитываем чётную степень:
\[
(x-2)^2\ge0.
\]

Поэтому
\[
(x-2)^2(x+7)\le0
\]
выполняется при
\[
x\le-7
\]
и отдельно при
\[
x=2.
\]

Получаем
\[
(-\infty;-7]\cup\{2\}.
\]

Пересечение двух множеств пусто:
\[
\boxed{\emptyset}.
\]

Если те же два условия образовали бы совокупность, потребовалось бы
объединить полученные множества.

% -------------------------------------------------

\topic{5. Текстовые задачи: таблица как математическая модель}

\key{Главный принцип.}
Сначала внесите искомую величину в подходящий столбец таблицы.
Остальные величины выражайте через основную формулу задачи.

\subsection*{Движение}

Основная связь:
\[
S=vt.
\]

Отсюда
\[
t=\frac{S}{v}.
\]

Если собственная скорость теплохода равна \(v\), скорость течения
равна \(x\), то
\[
v_{\text{по течению}}=v+x,
\]
\[
v_{\text{против течения}}=v-x.
\]

Рассмотрим ситуацию:

\begin{itemize}
  \item расстояние в каждую сторону --- \(285\) км;
  \item собственная скорость --- \(34\) км/ч;
  \item стоянка --- \(19\) ч;
  \item продолжительность всего путешествия --- \(36\) ч;
  \item скорость течения --- \(x\) км/ч.
\end{itemize}

Таблица:

\begin{table}[ht]
\centering
\caption{Движение теплохода}
\begin{tabularx}{\linewidth}{@{}Xccc@{}}
\toprule
Этап & Скорость, км/ч & Путь, км & Время, ч \\
\midrule
По течению
&
\(34+x\)
&
\(285\)
&
\(\dfrac{285}{34+x}\)
\\[0.6em]

Стоянка
&
\(0\)
&
\(0\)
&
\(19\)
\\

Против течения
&
\(34-x\)
&
\(285\)
&
\(\dfrac{285}{34-x}\)
\\[0.6em]
\bottomrule
\end{tabularx}
\end{table}

Получаем уравнение:
\[
\frac{285}{34+x}
+
19
+
\frac{285}{34-x}
=
36.
\]

Физические ограничения:
\[
0<x<34.
\]

Ограничение \(x<34\) гарантирует положительную скорость движения
против течения.

\subsection*{Производительность}

Если весь объём работы принят за \(1\), при времени выполнения \(t\)
производительность равна
\[
p=\frac1t.
\]

При совместной работе производительности складываются:
\[
p_1+p_2=p_{\text{совм}}.
\]

Пусть:

\begin{itemize}
  \item первая труба заполняет бассейн за \(19\) ч;
  \item вторая труба --- за \(x\) ч;
  \item вместе трубы заполняют бассейн за \(57\) мин.
\end{itemize}

Переводим минуты в часы:
\[
57\text{ мин}=\frac{57}{60}\text{ ч}.
\]

Тогда
\[
\frac1{19}
+
\frac1x
=
\frac1{57/60}.
\]

Удобно убрать трёхэтажную дробь:
\[
\frac1{19}
+
\frac1x
=
\frac{60}{57}.
\]

При этом
\[
x>0.
\]

\subsection*{Растворы и смеси}

Используйте три величины:

\begin{itemize}
  \item масса раствора;
  \item концентрация;
  \item масса растворённого вещества.
\end{itemize}

Основная формула:
\[
m_{\text{вещества}}
=
m_{\text{раствора}}\cdot c,
\]
где \(c\) --- концентрация в виде десятичной дроби.

Если смешаны равные массы \(m\) растворов концентраций \(10\%\)
и \(12\%\), итоговая масса равна
\[
2m.
\]

Пусть итоговая концентрация равна \(x\). Тогда
\[
0{,}10m+0{,}12m=2mx.
\]

Следовательно,
\[
0{,}22m=2mx.
\]

При \(m>0\):
\[
x=\frac{0{,}22}{2}=0{,}11.
\]

Переводим результат в проценты:
\[
\boxed{11\%}.
\]

\key{Проверка здравого смысла.}
При смешивании растворов одного вещества без дополнительных
процессов итоговая концентрация должна лежать между концентрациями
исходных растворов.

% -------------------------------------------------

\topic{6. Кусочно заданная функция и горизонталь \(y=m\)}

Рассмотрим функцию
\[
f(x)=
\begin{cases}
2{,}5x-3{,}5, & x<2,\\
-3x+7{,}5, & 2\le x\le3,\\
x-4{,}5, & x>3.
\end{cases}
\]

Каждая ветвь является линейной функцией. Для построения каждой ветви
достаточно двух точек.

Особенно полезно использовать граничные значения
\[
x=2
\qquad\text{и}\qquad
x=3.
\]

\begin{table}[ht]
\centering
\caption{Опорные точки ветвей}
\begin{tabular}{@{}ccc@{}}
\toprule
Ветвь & Значения \(x\) & Значения \(y\) \\
\midrule
\(y_1=2{,}5x-3{,}5\)
&
\(0,\;2\)
&
\(-3{,}5,\;1{,}5\)
\\

\(y_2=-3x+7{,}5\)
&
\(2,\;3\)
&
\(1{,}5,\;-1{,}5\)
\\

\(y_3=x-4{,}5\)
&
\(3,\;4\)
&
\(-1{,}5,\;-0{,}5\)
\\
\bottomrule
\end{tabular}
\end{table}

При \(x=2\) первая формула действует только для \(x<2\).
Вторая ветвь содержит значение \(x=2\), поэтому точка
\[
(2;1{,}5)
\]
принадлежит итоговому графику.

Аналогично точка
\[
(3;-1{,}5)
\]
принадлежит второй ветви.

\begin{center}
\begin{tikzpicture}

\begin{axis}[
  width=0.92\linewidth,
  height=9cm,
  axis lines=middle,
  unit vector ratio=1 1,
  xmin=-1.5,
  xmax=7,
  ymin=-5,
  ymax=4,
  samples=220,
  clip=false,
  grid=both,
  minor tick num=1,
  xlabel={$x$},
  ylabel={$y$},
  every axis x label/.style={
    at={(current axis.right of origin)},
    anchor=west
  },
  every axis y label/.style={
    at={(current axis.above origin)},
    anchor=south
  }
]

\addplot[
  accent,
  very thick,
  domain=-1.5:2
]
{2.5*x-3.5};

\addplot[
  darkaccent,
  very thick,
  domain=2:3
]
{-3*x+7.5};

\addplot[
  accent,
  very thick,
  domain=3:7
]
{x-4.5};

\addplot[
  only marks,
  mark=*,
  mark size=2.5pt,
  darkaccent
]
coordinates {
  (2,1.5)
  (3,-1.5)
};

\addplot[
  muted,
  dashed,
  thick,
  domain=-1.2:6.7
]
{1.5};

\addplot[
  muted,
  dashed,
  thick,
  domain=-1.2:6.7
]
{-1.5};

\node[
  muted,
  anchor=south west
]
at (axis cs:5.5,1.5)
{$y=1{,}5$};

\node[
  muted,
  anchor=north west
]
at (axis cs:5.5,-1.5)
{$y=-1{,}5$};

\end{axis}
\end{tikzpicture}
\end{center}

\subsection*{Как работать с параметром \(m\)}

Уравнение
\[
y=m
\]
задаёт горизонтальную прямую.

Чтобы определить число общих точек с графиком функции:

\begin{enumerate}
  \item мысленно перемещайте горизонтальную прямую вверх и вниз;
  \item считайте точки пересечения;
  \item отдельно проверяйте уровни, проходящие через граничные точки
  ветвей.
\end{enumerate}

При
\[
m=1{,}5
\]
получаются ровно две общие точки.

При
\[
m=-1{,}5
\]
также получаются ровно две общие точки.

Следовательно,
\[
\boxed{m=-1{,}5\quad\text{или}\quad m=1{,}5}.
\]

На экзамене полезно провести обе подходящие горизонтали на самом
чертежe.

% -------------------------------------------------

\topic{7. Типичные ошибки и финальная самопроверка}

\begin{itemize}

  \item
  \key{ОДЗ.}
  Знаменатель всегда отличен от нуля. Ограничения, появившиеся до
  сокращения, сохраняются до конца решения.

  \item
  \key{Строгость неравенства.}
  При знаках \(>\) и \(<\) нули числителя открыты.
  При \(\ge\) и \(\le\) нули числителя могут входить в ответ.

  \item
  \key{Нули знаменателя.}
  Все значения, обращающие знаменатель в ноль, исключаются.

  \item
  \key{Кратность корня.}
  При чётной кратности знак на соседних интервалах сохраняется.

  \item
  \key{Перенос членов.}
  Следите за изменением знака каждого переносимого слагаемого.

  \item
  \key{Домножение неравенства.}
  Перед умножением на выражение с переменной требуется знать его знак.

  \item
  \key{Система и совокупность.}
  Для системы используется пересечение, для совокупности --- объединение.

  \item
  \key{Единицы измерения.}
  Минуты и часы приводятся к одной единице до составления уравнения.

  \item
  \key{Производительность.}
  При совместной работе складываются производительности.

  \item
  \key{Растворы.}
  Проценты удобно переводить в десятичные дроби до вычисления массы
  растворённого вещества.

  \item
  \key{Кусочная функция.}
  Каждая формула используется только в указанной области.

  \item
  \key{Граничные точки графика.}
  Проверяйте принадлежность точки всей функции с учётом соседних ветвей.

  \item
  \key{Параметр \(m\).}
  В задаче с горизонталью \(y=m\) число пересечений должно быть видно
  непосредственно на графике.

\end{itemize}

\textbf{Приоритет следующего повторения:}
графики функций с модулем. Хорошей базой для этой темы служат уверенное
чтение кусочных функций, работа с граничными точками и преобразования
графиков.

% =================================================
% ТРЕНИРОВОЧНЫЙ БЛОК
% =================================================

\clearpage

\topic{Тренировочный блок}

\textbf{10 упражнений: 3 средних, 6 выше среднего, 1 сложное.}

Решения оформляйте последовательно. Для рациональных выражений
обязательно указывайте ОДЗ.

\begin{enumerate}[label=\textbf{\arabic*.}]

\item
\textbf{Средняя.}
Решите уравнение
\[
\frac{x^2}{4}+3x+8=0.
\]

\item
\textbf{Средняя.}
Решите неравенство
\[
(x-3)(x-3-\sqrt5)<0.
\]

\item
\textbf{Средняя.}
Решите систему
\[
\begin{cases}
3x^2+y=7,\\
4x^2-y=0.
\end{cases}
\]

\item
\textbf{Выше среднего.}
Решите неравенство
\[
\frac1x-\frac1{x-3}>0.
\]

\item
\textbf{Выше среднего.}
Решите уравнение
\[
(x-1)(x+2)^2-4(x+2)=0.
\]

\item
\textbf{Выше среднего.}
Решите систему неравенств
\[
\begin{cases}
\displaystyle
\frac{(x+2)(x-3)}{x+5}\ge0,
\\[0.6em]
(x-2)^2(x+7)\le0.
\end{cases}
\]

\item
\textbf{Выше среднего.}
Теплоход проходит \(285\) км по течению реки, делает стоянку
продолжительностью \(19\) ч и возвращается обратно.
Собственная скорость теплохода равна \(34\) км/ч.
Всё путешествие длится \(36\) ч.
Найдите скорость течения.

\item
\textbf{Выше среднего.}
Первая труба заполняет бассейн за \(19\) ч.
Две трубы вместе заполняют его за \(57\) мин.
За сколько часов заполнит бассейн вторая труба, работая одна?

\item
\textbf{Выше среднего.}
Смешали \(10\%\)-й и \(16\%\)-й растворы одного вещества и получили
\(18\) кг раствора концентрации \(14\%\).
Найдите массу каждого исходного раствора.

\item
\textbf{Сложная.}
Постройте график функции
\[
f(x)=
\begin{cases}
2{,}5x-3{,}5, & x<2,\\
-3x+7{,}5, & 2\le x\le3,\\
x-4{,}5, & x>3,
\end{cases}
\]
и найдите все значения \(m\), при которых прямая \(y=m\) имеет
с графиком ровно две общие точки.

\end{enumerate}

% =================================================
% ОТВЕТЫ
% =================================================

\clearpage

\topic{Ответы к тренировочному блоку}

\begin{table}[ht]
\centering
\caption{Краткие ответы}
\begin{tabularx}{\linewidth}{@{}cX@{}}
\toprule
\textbf{№} & \textbf{Ответ} \\
\midrule

1
&
\(x=-8,\,-4\).
\\

2
&
\(x\in(3;\,3+\sqrt5)\).
\\

3
&
\((1;4)\), \((-1;4)\).
\\

4
&
\(x\in(0;3)\).
\\

5
&
\(x=-3,\,-2,\,2\).
\\

6
&
\(\emptyset\).
\\

7
&
\(4\) км/ч.
\\

8
&
\(1\) ч.
\\

9
&
\(6\) кг \(10\%\)-го раствора и \(12\) кг \(16\%\)-го раствора.
\\

10
&
\(m=-1{,}5\) или \(m=1{,}5\).
\\

\bottomrule
\end{tabularx}
\end{table}

\vspace{1em}

\textbf{Финальная проверка перед сдачей:}

\begin{itemize}

  \item
  в дробных выражениях указаны все ограничения;

  \item
  в неравенствах корректно оформлены открытые и закрытые точки;

  \item
  нули знаменателей исключены;

  \item
  в системах найдено пересечение решений;

  \item
  в текстовых задачах единицы измерения согласованы;

  \item
  в задачах на движение проверены физические ограничения;

  \item
  на графике отмечены граничные точки ветвей;

  \item
  для задачи с \(y=m\) проведены горизонтали, подтверждающие ответ.

\end{itemize}

\end{document}